\chapteropen{04}

\begin{cele}
\item policzyć dzień pracy razem z przerwami, dojazdem i powrotem, zamiast pytać, ile godzin „zajmuje praca”;
\item rozpoznać, który element dnia pracy jest kosztem, a który czasem na regenerację;
\item przygotować argumenty do rozmowy o pracy zdalnej, elastycznym początku dnia lub krótszym dniu;
\item ustawić granicę wieczorną, która działa także wtedy, gdy powiadomienia nie przestają przychodzić;
\item ocenić uczciwie, kiedy plan dnia wokół 12 godzin snu przestaje działać, czyli przy pracy zmianowej i nocnej.
\end{cele}

\section{Osiem godzin, które nie są ośmioma}

Umowa mówi o ośmiu godzinach pracy. Doba mówi o dwudziestu czterech godzinach i to ona zostaje w budżecie. Zanim zaczniesz cokolwiek przesuwać, policz wszystko, co dzień pracy zabiera, a nie tylko czas przy biurku lub przy laptopie. Osiem godzin pracy to rzadko osiem godzin z dnia.

Weźmy typowy dzień biurowy. Wychodzisz z domu, dojeżdżasz, siedzisz osiem godzin, wracasz. Do tego dochodzi przerwa na obiad, kilka krótkich przerw, dojście do kuchni po kawę i ten kwadrans, w którym przeglądasz pocztę jeszcze przed włączeniem komputera. Jeżeli część z tego jest wliczona w osiem godzin, a część nie, to dokładna liczba zależy od regulaminu i od tego, jak twój pracodawca liczy czas. Dlatego w tym rozdziale nie zakładamy żadnej konkretnej polityki firmy, tylko liczymy to, co zawsze trzeba policzyć: ile godzin doby zajmuje praca i jej otoczenie.

Dzień biurowy z dojazdem w budżecie tej książki wygląda tak. Szablon B z rozdziału szóstego to plan zbudowany wokół jednej reguły: sen zajmuje 12 godzin, praca 8, a reszta musi się zmieścić w tym, co zostaje.

\generated{week-templates}

W szablonie B dojazd to jedna godzina, czyli przy 8 godzinach pracy dzień zajmuje już dziewięć godzin tylko na pracę i dojazd razem. Wolnego zostaje \val{week.B.free} godziny, czyli piętnaście minut. Bez dojazdu, jak w szablonie A, pracy zdalnej, wolnego jest \val{week.A.free} godziny, czyli trzydzieści minut. Różnica między tymi dwoma szablonami wynosi więc zaledwie kwadrans wolnego, a dojazd zabiera godzinę: arytmetyka jest prosta, a wnioski z niej nie zawsze są wygodne.

Osiem godzin pracy nie jest stałą, którą można dowolnie poszerzać lub skracać. Jest to umowna wielkość, którą pracodawca i pracownik ustalają, a w wielu krajach prawo wyznacza jej górną granicę i minimalny odpoczynek między dniami pracy.\verify{Polskie przepisy o czasie pracy (Kodeks pracy) — wymiar 8 godzin dziennie i minimalny dobowy odpoczynek; sprawdzić aktualne brzmienie artykułów przed powołaniem się na nie w rozmowie z pracodawcą.} To ważne dla dwóch rzeczy: po pierwsze, granica dnia nie jest twoim wymysłem, więc możesz się na nią powołać; po drugie, nadgodziny nie są czymś, co „po prostu się zdarza”, tylko odstępstwem, które ktoś musi zaakceptować.

\section{Dojazd jako koszt, nie jako przerwa}

Dojazd bywa traktowany jak czas wolny, bo nie siedzisz przy biurku. Z punktu widzenia snu to nie jest czas wolny. Słuchanie podcastu w autobusie nie przywraca sił, a stanie w korku nie jest regeneracją. Jest to godzina, w której nie śpisz, nie gotujesz, nie ćwiczysz i nie odpoczywasz, a jednocześnie jesteś zmęczony i zdekoncentrowany.

Policz dojazd tygodniowo, a nie dziennie. Godzina dziennie przy pięciu dniach pracy daje pięć godzin tygodniowo. Dla porównania, wolny czas w szablonie A to 0,5 godziny dziennie, czyli około trzech i pół godziny w tygodniu. Dwie godziny dziennie to dziesięć godzin tygodniowo, czyli ponad dwie piąte doby. Ta liczba zwykle zmienia decyzję bardziej niż każda rozmowa o motywacji.

Nie każdy dojazd jest równie drogi. Samochód do pracy odbiera czas, ale daje kontrolę nad godziną wyjazdu. Komunikacja publiczna pozwala czytać, spać w drodze lub nie robić nic, za to zależy od rozkładu, przesiadek i opóźnień. Najgorszy wariant to dojazd, którego długość zmienia się w zależności od dnia tygodnia. Wtedy nie wiesz, czy dzień ma osiem godzin, czy dziewięć, i plan staje się fikcją.

\begin{falle}
Częsty błąd: „mieszkam blisko, więc dojazd praktycznie nie istnieje”. Liczy się czas od drzwi do drzwi, z przebraniem się, wejściem do budynku i windą. Zmierz to przez tydzień zwykłym stoperem, zanim uznasz dojazd za zerowy. Pięć minut w jedną stronę to dwadzieścia minut dziennie, czyli około dwóch godzin tygodniowo.
\end{falle}

\section{Negocjowanie: co naprawdę można zaproponować}

Pracodawcy rzadko zgadzają się na zmianę pracy wokół snu, bo sen nie jest argumentem, który ktokolwiek w firmie rozumie. Argumentem jest wynik pracy, dostępność i koszt. Dlatego rozmowa o pracy zdalnej, elastycznych godzinach albo późniejszym początku dnia powinna zaczynać się od tego, co zespół dostaje w zamian, a nie od tego, czego potrzebujesz.

Są trzy rodzaje prośby, które mają szansę przejść, i każdy ma inny koszt dla firmy.

Pierwszy to praca zdalna przez część tygodnia. Najłatwiej przyjąć dzień bez spotkań, w którym zadania wymagają skupienia. Propozycja powinna określać, które dni, jaka jest dostępność w godzinach wspólnych i jak mierzysz wynik. Jeżeli zespół nie ma wspólnych godzin, zaproponuj je: dwie godziny rano w dni zdalne, na przykład między dziesiątą a dwunastą, pokrywają większość pytań.

Drugi to elastyczny początek i koniec dnia. Tu najczęściej pojawia się opór, bo elastyczność bywa rozumiana jako brak zasad. Dlatego zaproponuj konkretny przedział, na przykład start między siódmą a dziewiątą trzydzieści, z zachowaniem ośmiu godzin pracy i z jednym stałym oknem, w którym cały zespół jest dostępny.

Trzeci to skrócenie dnia, na przykład do siedmiu godzin i czterdziestu pięciu minut, z odpowiednim obniżeniem zakresu. Jest to najtrudniejsza prośba, bo wymaga rozmowy o pensji albo o zadaniach, ale też najbardziej uczciwa wobec tego, co faktycznie dzieje się w dniu pracy.

Przygotuj rozmowę w trzech krokach. Opisz, co chcesz zmienić, w jednym zdaniu. Pokaż, jak zmiana wpływa na zespół, najlepiej na konkretnych zadaniach lub terminach. Zaproponuj okres próbny, na przykład cztery tygodnie, z ustalonym momentem przeglądu. Okres próbny zmniejsza koszt decyzji dla obu stron i daje dane zamiast przekonań.

Nie obiecuj efektów zdrowotnych ani nie tłumacz się snem w sposób, który brzmi jak prośba o ulgę. Mówienie „potrzebuję stałej pory” jest prawdziwe, ale nie jest argumentem. Argumentem jest „przy stałym starcie o ósmej rano dowożę raporty w terminie, a od kiedy przesunąłem start, liczba poprawek spadła o połowę”, jeżeli to rzeczywiście jest prawda. Jeżeli tego nie możesz zmierzyć, nie wstawiaj liczby.

\begin{przyklad}
Anna pracuje osiem godzin i dojeżdża 70 minut w jedną stronę, czyli 140 minut dziennie. Plan mówi, że dzień ma 24 godziny, sen 12, a praca z dojazdem zabiera 8 godzin i 140 minut, czyli 10 godzin i 20 minut. Zostaje 1 godzina i 40 minut na wszystko poza snem i pracą: jedzenie, higienę, ruch, dom i wolny czas. To mało. Anna proponuje pracę zdalną w dwa dni w tygodniu. Przy tej zmianie tygodniowy dojazd spada z 700 minut do 420 (trzy dni dojazdu po 140 minut), czyli o 280 minut, prawie pięć godzin. Dwa dni w tygodniu nie zmieniają dnia pracy, ale zmieniają budżet tygodnia. Żadna z tych liczb nie mówi, czy Anna będzie lepiej spała. Mówi tylko, ile czasu zyskuje.
\end{przyklad}

\section{Asynchroniczność i okno odpowiedzi}

Praca asynchroniczna wygląda na łatwiejszą, bo nie wymaga bycia dostępnym w każdej chwili. W praktyce często jest trudniejsza, bo zamienia jedną ciągłą presję na wiele krótkich. Wiadomość w komunikatorze, mail z pytaniem, powiadomienie w systemie zgłoszeń: każde z nich jest drobne, ale razem tworzą dzień, w którym nigdy nie jesteś naprawdę w jednym zadaniu.

Rozwiązaniem nie jest ignorowanie wiadomości, tylko ustalenie okna odpowiedzi. Jeżeli wiadomości zespołu czekają najwyżej cztery godziny, a ty sprawdzasz skrzynkę w dwóch stałych momentach dnia, zespół wie, kiedy dostanie odpowiedź, i nie musi zgadywać. To jest umowa, którą można zaproponować. Zmiana „zawsze odpowiadam od razu” na „odpowiadam do końca dnia, a pilne sprawy oznaczam” zwykle przechodzi łatwiej, niż się wydaje.

Warto też rozróżnić, co jest pilne, a co tylko głośne. Pilne jest to, co blokuje czyjąś pracę do końca dnia, a tego w większości zespołów jest niewiele. Reszta może poczekać do następnego okna. Zapisanie tej reguły na jednej stronie, z przykładami, zmniejsza liczbę sytuacji, w których trzeba negocjować każdą wiadomość osobno.

\section{Granica wieczorna}

Najtrudniejszy element dnia pracy nie leży w godzinach pracy. Leży w tym, co dzieje się po nich. Wieczorem zerkasz na telefon, potem jeszcze raz przed snem, a skrzynka w tle dopisuje sobie kolejne pozycje. Jeżeli wieczór nie ma granicy, dzień pracy wydłuża się o czas, którego nikt nie wpisał w umowę, a sen zaczyna się później, niż planowałeś.

Granica ma trzy części. Pierwsza to wyłączenie powiadomień z aplikacji służbowych w wybranej godzinie, na przykład od siedemnastej trzydzieści, i zostawienie ich dla jednego kanału, jeżeli firma wymaga dyżuru. Druga to fizyczne oddzielenie urządzenia: telefon służbowy poza sypialnią, a najlepiej poza domem w czasie wolnym. Trzecia to rytuał zamknięcia dnia.

Rytuał zamknięcia to kilka minut na koniec pracy: spisanie trzech spraw na jutro, zapisanie jednej rzeczy, która została w połowie, i zamknięcie wszystkich otwartych kart. Jego zadaniem nie jest zwiększanie produktywności. Ma przekazać mózgowi informację, że sprawy są zapisane, więc nie trzeba ich pamiętać. Ten efekt jest znany z badań nad pamięcią zadań, ale nie opieramy na nim żadnej obietnicy.\verify{Efekt domknięcia zadań i zapisywania planu na następny dzień (w literaturze nazywany efektem Zeigarnik) — przed powołaniem się na jego wielkość sprawdzić oryginalne badania i replikacje; w tym rozdziale jest to wskazówka praktyczna, nie zmierzony wynik.}

Granica działa tylko wtedy, gdy ma koszt dla tego, kto ją łamie. Jeżeli w firmie od dyżuru oczekuje się odpowiedzi w nocy, granica nie jest dobrowolna. Wtedy trzeba ją wynegocjować jako dyżur z dodatkowym odpoczynkiem albo odmówić, a nie udawać, że nie ma problemu.

\begin{zasada}
Wieczór ma koniec, który jest w kalendarzu, a nie w głowie. Wpisz go tak samo jak spotkanie: „koniec pracy 17:30, zamknięcie dnia 17:30–17:40, wyłączenie powiadomień 17:40”. Jeżeli nie da się wpisać, nie da się też dotrzymać.
\end{zasada}

\section{Rozmowa o dostępności z szefem i zespołem}

Dostępność to temat, który w wielu zespołach nigdy nie został nazwany. Każdy zakłada, że inni są dostępni cały czas, i nikt nie wie, co naprawdę oznacza „dostępny”. Zanim poprosisz o zmianę, nazwij to, co jest. Jeżeli oczekiwania są rozmyte, każda zmiana wygląda jak wyjątek, a wyjątki są podejrzane.

Porozmawiaj z przełożonym o trzech rzeczach: w jakich godzinach zespół musi być osiągalny, jak szybko oczekuje odpowiedzi na sprawy zwykłe i pilne, oraz co się dzieje, gdy ktoś nie jest dostępny. Ostatni punkt jest najważniejszy, bo pokazuje, czy firma ma zapas na nieobecność, czy opiera się na jednej osobie.

Uważaj na pułapkę „zawsze jestem, tylko czasem nie odpisuję”. Ludzie, którzy tak mówią, zwykle nie mają w kalendarzu żadnego bloku, w którym nikt ich nie zapyta. To nie jest dostępność, tylko brak granic, który zostanie przypisany tobie.

\section{Praca zmianowa i nocna: gdzie plan się nie sprawdzi}

Ta książka buduje dzień wokół 12 godzin snu i stałej pory. Praca zmianowa i nocna sprawiają, że stała pora przestaje być dostępna, bo zmiana przesuwa godzinę, o której kładziesz się i wstajesz. Nie ma tu sprytnego triku, który sprawi, że plan zadziała tak samo. Trzeba powiedzieć to wprost.

Przy pracy nocnej cykl dnia i nocy rozjeżdża się z godzinami, w których pracujesz i śpisz. Sen w dzień jest zwykle krótszy i płytszy, a światło słoneczne, hałas i obowiązki domowe przeszkadzają w nim. Dlatego 12 godzin snu w ciągu doby nocnej jest znacznie trudniejsze do uzyskania niż w ciągu nocy zwykłej.\verify{Jakość i długość snu w pracy zmianowej i nocnej w porównaniu z pracą dzienną — przeglądy badań nad pracą zmianową; konkretne wartości zależą od rodzaju zmiany i nie są tu podawane.}

Co z tego wynika praktycznie? Po pierwsze, nie próbuj utrzymać tej samej stałej pory przy każdej zmianie. Lepiej zaplanować dwa lub trzy stałe rytmy i przechodzić między nimi świadomie, z kilkoma dniami przejścia. Po drugie, zaciemnij sypialnię i ustal, co dzieje się z dziećmi, zakupami i domem w godzinach, w których śpisz. Po trzecie, zaakceptuj, że „dzień wokół snu” w tym przypadku oznacza dzień wokół snu zmiennego, a szablony z tej książki służą jako punkt wyjścia do zmiany, a nie gotowy plan.

Jest też granica, której nie należy przekraczać: jeżeli sen na zmianie systematycznie nie wystarcza, a objawy się utrzymują, to nie jest problem planowania, tylko sygnał do rozmowy z lekarzem. Ten rozdział nie zastępuje takiej oceny.

\section{Nadgodziny jako zjadacz snu}

Nadgodziny rzadko wyglądają jak zagrożenie. Jedna godzina dziś, jedna jutro, „zostanę, żeby to domknąć”. Każda z nich osobno wygląda niewinnie. Razem zamieniają się w wydłużony dzień, który musi skądś zabrać czas, a najłatwiej zabrać go ze snu, czyli z bloku, który wydaje się najbardziej elastyczny.

Policzmy to dla szablonu A. Ma on \val{week.A.free} godziny wolnego w ciągu doby. Dzień z jedną godziną nadgodzin przekracza więc ten zapas o pół godziny: wolne, które miało być buforem, znika, i trzeba je odebrać z innego bloku. Najczęściej jest to posiłek, ruch albo sen, czyli trzy bloki, które już są ustawione na minimum. Jeżeli nadgodziny zdarzają się codziennie przez tydzień, różnica wynosi pięć godzin i nie da się jej odzyskać z wolnego czasu w weekend bez kosztu dla regeneracji.

Dlatego nadgodziny trzeba traktować jak decyzję, a nie jak nieuniknioną część pracy. Każda nadgodzina powinna mieć nazwę: co konkretnie została zrobiona, dla kogo i czy da się ją przesunąć na jutro. Jeżeli odpowiedź brzmi „nie da się”, to pytanie należy do przełożonego, a nie do twojego snu.

\begin{uwaga}
Odbiór nadgodzin w czasie wolnym też ma koszt. Jeżeli odbierasz dwie godziny w piątek po południu, a w poniedziałek znowu zaczynasz o siódmej, to nie zrównoważyłeś tygodnia, tylko przesunąłeś dług. Liczy się tydzień, nie pojedynczy dzień.
\end{uwaga}

\begin{cwiczenia}
\exitem{Ile godzin doby zajmuje praca, jeżeli dzień to osiem godzin pracy, godzina dojazdu w jedną stronę i powrót tym samym czasem?}{Dojazd w obie strony to dwie godziny, więc razem 8 + 2 = 10 godzin. Dzień zostawia wtedy 24 − 12 − 10 = 2 godziny na wszystko poza snem i pracą.}
\exitem{Dlaczego dojazd tygodniowo, a nie dziennie, zmienia decyzję o pracy zdalnej?}{Dzienna godzina wygląda mało, ale przy pięciu dniach pracy daje pięć godzin tygodniowo. Wolny czas szablonu A, 0,5 godziny dziennie, daje w tym samym tygodniu około trzech i pół godziny. Tygodniowa suma pokazuje skalę kosztu.}
\exitem{Porównaj szablony A i B z tabeli. Ile wolnego zostaje w każdym i jaka jest różnica?}{Szablon A ma 0,5 h wolnego, szablon B 0,25 h. Różnica wynosi 0,25 h, czyli 15 minut, a dojazd w szablonie B zabiera 1 godzinę.}
\exitem{Jakie trzy elementy powinna zawierać rozmowa o pracy zdalnej, żeby miała szansę przejść?}{Opis zmiany w jednym zdaniu, pokazanie wpływu na zespół na konkretnych zadaniach lub terminach oraz propozycja okresu próbnego z ustalonym przeglądem.}
\exitem{Co zrobić, gdy wiadomości zespołu wymagają odpowiedzi w ciągu minut, ale ty chcesz pracować bez przerwania?}{Ustalić okno odpowiedzi, na przykład cztery godziny, i określić, co jest pilne. Pilne sprawy mają osobny kanał, a reszta czeka do najbliższego okna. Zespół dostaje to na piśmie.}
\exitem{Dzień ma nadgodzinę ponad osiem godzin pracy. Przy szablonie A wolnego jest 0,5 h. Ile godzin trzeba odebrać z innego bloku i z jakiego?}{Nadgodzina zabiera 1 h, wolnego jest 0,5 h, więc brakuje 0,5 h. Brakującą część najczęściej odbiera się z posiłków, ruchu lub snu, czyli z bloków już ustawionych na minimum.}
\exitem{Dlaczego granica wieczorna wymaga kosztu dla łamiącego ją?}{Bo bez kosztu granica jest dobrowolna, a wszystko, co dobrowolne, ustępuje pierwszej pilnej wiadomości. Jeżeli firma oczekuje odpowiedzi w nocy, trzeba to wynegocjować jako dyżur z dodatkowym odpoczynkiem.}
\end{cwiczenia}

\zadanie{Przez pięć dni roboczych zmierz dojazd od drzwi do drzwi z przebraniem i wejściem, a także czas faktycznej pracy według wpisów w kalendarzu. Zapisz oba czasy w tabeli. Na koniec tygodnia policz sumę dojazdu i porównaj ją z tym, co zakładasz w swoim szablonie. Miara sukcesu: masz pięć zmierzonych dni i jedną liczbę tygodniowego dojazdu, którą możesz przedstawić przełożonemu.}
